\scanpage{215}
\mspage{131}
% archive sheet number: 214
% element: diagram | top, first two lines | components to x alpha and x beta
\diagram{\begin{tikzcd}[ampersand replacement=\&]
x_{\alpha,\beta,\unsure{\sigma|\gamma}}\arrow[r,Rightarrow,"\unsure{v|\nu}_{\alpha,\beta,\sigma}"]\arrow[dr,dash,"\unsure{w|\omega}_{\alpha,\beta,\sigma}"']\&x_\alpha\\
\&x_\beta,
\end{tikzcd}}
% element: text | upper third, lines 3-9 | donc pour
donc pour $x_{\alpha,\beta,\sigma}\in\mathrm{Ob}\,C_0$ fixés, les choix possibles pour $u_{\alpha,\beta,\sigma}$ sont de cardinal majoré par $\pi_1\cdot\pi_1^{\unsure{2|3}}=\pi_1$, de plus on a $\operatorname{card}\mathrm{Ob}\,C_0\leq\pi_1$ donc le nb \deleted{\unsure{total}} des flèches $u_{\alpha,\beta,\sigma}$ \deleted{$\unsure{v_{\alpha,\beta,\sigma},w_{\alpha,\beta,\sigma}}$ \unsure{est}} \inserted{(pour $\alpha,\beta$ fixés)} est majoré par $\pi_1\cdot\pi_1=\pi_1$. Donc le nb. total des flèches $u_{\alpha,\beta,\sigma}$, pour $\alpha,\beta,\sigma$ variables, est majoré par $\pi\cdot\pi_1=\pi$ (si $\pi\geq\pi_1$).

% element: text | middle, lines 10-12 | Donc x est représenté
Donc $x$ est représenté par une catégorie $I$, ayant \deleted{\unsure{les mêmes}} \inserted{comme} objets \deleted{que $B$} \inserted{\unsure{les élém. de}} \deleted{\unsure{la}} somme de $A$ et de
% element: deletion | middle right, overwritten expression | indexed coproduct
\deleted{$\unsure{\coprod_{\alpha,\beta}\Sigma_{\alpha,\beta}|p_{\alpha,\beta}\Sigma_{\alpha,\beta}|\coprod_{\alpha,\beta}\{\sigma\}}$}
% element: deletion | middle, crossed three-line former description | mais les flèches non identiques
\deleted{\unsure{mais} les flèches non identiques de $B$ \unsure{ét} \unsure{étant|A. et} \unsure{remplacées chacune} par les paquets de flèches $v_{\alpha,\beta,\unsure{\lambda|\sigma}},w_{\alpha,\beta,\lambda}$}
% element: text | lower third, lines 15-19 | des indices triples
des indices triples $\alpha,\beta,\sigma$, et pour chaque tel indice, existent deux flèches $v_{\alpha,\beta,\sigma},w_{\alpha,\beta,\sigma}$ non identiques de source cet \deleted{\unsure{indice}} objet, \deleted{\unsure{les|en}} et de but respectivement en $\alpha,\beta$. On a \deleted{\unsure{en fait}} donc $\operatorname{card}\mathrm{Fl}\,I\leq\pi$, et on a un foncteur
% element: diagram | lower third, line 20 | I to C
\diagram{\begin{tikzcd}[ampersand replacement=\&]
I\arrow[r]\&C
\end{tikzcd}}
% element: text | foot, lines 21-24 | alpha maps to x alpha
$\unsure{\alpha_\beta|\alpha}\unsure{\mathrel{\vert\!-\!\vert}|\mapsto}x_\alpha$, $(\alpha,\beta,\sigma)\mapsto x_{\alpha,\beta,\sigma}$, tel que $x=\lim_{i\in I}\unsure{x'|x^i|x_i}$. Cela achève la dém. \unsure{du|de} 1\textsuperscript{o}).
